\section{Consensus Taxonomy Summary}

\subsection{Winner by $k$}
\begin{table}[t]
\centering
\begin{tabular}{lrrr}
\hline
Slice & consensus\_all & sum\_minus\_emb & sum\_minus\_cc \\
\hline
top-$k$=6 & 66 & 30 & 0 \\
top-$k$=30 & 35 & 0 & 13 \\
Total & 101 & 30 & 13 \\
\hline
\end{tabular}
\caption{Bank-based local review winners by slice.}
\end{table}

\subsection{Taxonomy by winner}
\begin{table}[t]
\centering
\begin{tabular}{lrrrrr}
\hline
Taxonomy label & Count & Share & consensus\_all & sum\_minus\_emb & sum\_minus\_cc \\
\hline
broad\_context\_noise & 50 & 34.7\% & 47 & 2 & 1 \\
material\_family\_coherence & 37 & 25.7\% & 28 & 6 & 3 \\
method\_family\_coherence & 22 & 15.3\% & 14 & 8 & 0 \\
over\_regularized\_consensus & 16 & 11.1\% & 0 & 11 & 5 \\
application\_umbrella\_noise & 6 & 4.2\% & 3 & 2 & 1 \\
single\_cue\_specificity & 6 & 4.2\% & 2 & 1 & 3 \\
semantic\_drift & 4 & 2.8\% & 4 & 0 & 0 \\
coherent\_refinement & 3 & 2.1\% & 3 & 0 & 0 \\
\hline
\end{tabular}
\caption{Primary case taxonomy across all reviewed cases.}
\end{table}

\subsection{Representative cases}
\paragraph{consensus\_all}
\begin{itemize}
\item \textbf{Application of Poly(ethylene glycol)-based Aqueous Biphasic Systems as Reaction and Reactive Extraction Media} (field\_15\_k30\_sum\_minus\_cc\_vs\_consensus\_bank\_n48, $k$=30, broad\_context\_noise, gap=4.0). Winner advantage: Group A provides an excellent local neighborhood. The top-ranked papers focus specifically on the partitioning and phase diagrams of Poly(ethylene glycol)-based Aqueous Biphasic Systems (PEG-ABS), which is the exact system studied in the target paper.. Loser failure mode: Group B is almost entirely focused on a different class of solvents, Deep Eutectic Solvents (DES), and fails to capture the target's core topic of ABS, making it a poor representation of the immediate research context..
\item \textbf{Tunable thermomorphism and applications of ionic liquid analogues of Girard's reagents} (field\_15\_k30\_sum\_minus\_cc\_vs\_consensus\_bank\_n48, $k$=30, material\_family\_coherence, gap=4.0). Winner advantage: Group A provides an excellent and highly coherent neighborhood focused on the target's specific topic: using ionic liquids that exhibit temperature-dependent phase behavior (thermomorphism, LCST/UCST) for homogeneous liquid-liquid extraction.. Loser failure mode: Group B is a very broad and noisy collection of papers generally related to 'ionic liquids' but covering disparate applications like crystallization, electrolytes, and general synthesis, completely missing the target's specific context..
\item \textbf{Chemical Absorption of Sulfur Dioxide in Room-Temperature Ionic Liquids} (field\_15\_k30\_sum\_minus\_cc\_vs\_consensus\_bank\_n48, $k$=30, material\_family\_coherence, gap=4.0). Winner advantage: Group A provides a perfectly coherent local neighborhood where every paper is directly about the target's specific topic: SO2 absorption in ionic liquids. Group B is extremely weak; after the first paper, it diverges into papers about different gases (CO2, H2S) or the general properties of ionic liquids, completely missing the target's core research problem.. Loser failure mode: Group A provides a perfectly coherent local neighborhood where every paper is directly about the target's specific topic: SO2 absorption in ionic liquids. Group B is extremely weak; after the first paper, it diverges into papers about different gases (CO2, H2S) or the general properties of ionic liquids, completely missing the target's core research problem..
\item \textbf{A New Category of Liquid Salt--Liquid Ionic Phosphates (LIPs)} (field\_15\_k30\_sum\_minus\_cc\_vs\_consensus\_bank\_n48, $k$=30, material\_family\_coherence, gap=4.0). Winner advantage: Group A provides an exceptional local neighborhood. Its top-ranked paper is a direct continuation or companion paper to the target, sharing the key term "LIPs" (Liquid Ionic Phosphates) and focusing on polyammonium salts. The subsequent papers cohere tightly around the theme of dicationic/polycationic ionic liquids, which is the central structural motif of the target paper.. Loser failure mode: Group B is a very broad and unfocused collection of general reviews and disparate specific studies on ionic liquids, failing entirely to capture the target's specific contribution to the field..
\end{itemize}
\paragraph{sum\_minus\_cc}
\begin{itemize}
\item \textbf{Carbon Paste Electrode Modified with Functionalized Nanoporous Silica Gel as a New Sensor for Determination of Silver Ion} (field\_15\_k30\_sum\_minus\_cc\_vs\_consensus\_bank\_n48, $k$=30, single\_cue\_specificity, gap=3.0). Winner advantage: Group A provides an excellent and highly coherent neighborhood. Its top-ranked papers are directly relevant, focusing on either the same method (modified carbon paste electrodes) or the same target analyte (silver ion).. Loser failure mode: Group B is much broader and noisier, with many of its top-ranked papers focusing on the determination of organic drug molecules, which is a different application area from the target's focus on metal ion sensing..
\item \textbf{A single-site hydroxyapatite-bound zinc catalyst for highly efficient chemical fixation of carbon dioxide with epoxides} (field\_15\_k30\_sum\_minus\_cc\_vs\_consensus\_bank\_n48, $k$=30, single\_cue\_specificity, gap=3.0). Winner advantage: Group A provides an excellent local neighborhood by correctly identifying both the target's specific chemical reaction (CO2 fixation with epoxides) and the specific catalyst family (zinc-based catalysts).. Loser failure mode: Group B identifies the general reaction but completely misses the crucial zinc catalyst aspect, instead providing a diverse and less relevant set of catalysts (Ni, Ru, Cr, organocatalysts) and even including mismatched papers about polymerization..
\item \textbf{Preparation and Characterization of New Room Temperature Ionic Liquids} (field\_15\_k30\_sum\_minus\_cc\_vs\_consensus\_bank\_n48, $k$=30, material\_family\_coherence, gap=2.0). Winner advantage: Group A provides a more coherent and focused neighborhood. Its top-ranked papers are all centered on the synthesis and characterization of imidazolium-based ionic liquids, which directly matches the target's scope.. Loser failure mode: Group B is noisier, including papers on different cation types (guanidinium and pyridinium) that are less relevant to the target's specific chemistry, thereby diluting the immediate research context..
\item \textbf{Surface Analysis of Ionic Liquids with and without Lithium Salt Using X-ray Photoelectron Spectroscopy} (field\_15\_k30\_sum\_minus\_cc\_vs\_consensus\_bank\_n48, $k$=30, material\_family\_coherence, gap=2.0). Winner advantage: Group B provides a much more specific and relevant neighborhood for the target paper. It includes multiple papers studying the exact same ionic liquid ([EMIM][Tf2N]) and, critically, several papers investigating the interaction of this class of ionic liquids with lithium, which is the central perturbation in the target study. Group A is plausible, as it focuses on the correct technique (XPS) and general problem (ionic liquid surfaces with dissolved salts), but its focus on platinum and palladium complexes makes it less directly relevant than Group B's focus on lithium.. Loser failure mode: Group B provides a much more specific and relevant neighborhood for the target paper. It includes multiple papers studying the exact same ionic liquid ([EMIM][Tf2N]) and, critically, several papers investigating the interaction of this class of ionic liquids with lithium, which is the central perturbation in the target study. Group A is plausible, as it focuses on the correct technique (XPS) and general problem (ionic liquid surfaces with dissolved salts), but its focus on platinum and palladium complexes makes it less directly relevant than Group B's focus on lithium..
\end{itemize}
\paragraph{sum\_minus\_emb}
\begin{itemize}
\item \textbf{Impact of Ionic Liquid Physical Properties on Lipase Activity and Stability} (field\_15\_k06\_sum\_minus\_emb\_vs\_consensus\_bank\_n48, $k$=6, material\_family\_coherence, gap=3.0). Winner advantage: Group A provides a much better local neighborhood. Its top-ranked papers are directly about lipase or other enzymes in ionic liquids, which is the core topic of the target paper. Group B is extremely poor, with its top-ranked papers focusing on an entirely different application of ionic liquids (metal ion extraction), demonstrating a fundamental misunderstanding of the target's research context.. Loser failure mode: Group A provides a much better local neighborhood. Its top-ranked papers are directly about lipase or other enzymes in ionic liquids, which is the core topic of the target paper. Group B is extremely poor, with its top-ranked papers focusing on an entirely different application of ionic liquids (metal ion extraction), demonstrating a fundamental misunderstanding of the target's research context..
\item \textbf{Combinatorial Rheology of Branched Polymer Melts} (field\_15\_k06\_sum\_minus\_emb\_vs\_consensus\_bank\_n48, $k$=6, method\_family\_coherence, gap=3.0). Winner advantage: Group A provides an excellent, coherent neighborhood. Its top-ranked papers focus on the same core problem as the target: developing general algorithms and computational models to predict the linear rheology of complex branched polymers. The lower-ranked papers in A provide relevant context by focusing on specific architectures like comb and star polymers, which the target paper explicitly mentions. Group B is significantly weaker, as it buries the most relevant papers at the bottom of its ranking and includes a clear topical mismatch (rank 6, polymer film surface dynamics) which is distinct from the target's focus on bulk melts.. Loser failure mode: Group A provides an excellent, coherent neighborhood. Its top-ranked papers focus on the same core problem as the target: developing general algorithms and computational models to predict the linear rheology of complex branched polymers. The lower-ranked papers in A provide relevant context by focusing on specific architectures like comb and star polymers, which the target paper explicitly mentions. Group B is significantly weaker, as it buries the most relevant papers at the bottom of its ranking and includes a clear topical mismatch (rank 6, polymer film surface dynamics) which is distinct from the target's focus on bulk melts..
\item \textbf{Binuclear chromium–salan complex catalyzed alternating copolymerization of epoxides and cyclic anhydrides} (field\_15\_k06\_sum\_minus\_emb\_vs\_consensus\_bank\_n48, $k$=6, over\_regularized\_consensus, gap=3.0). Winner advantage: Group A provides an excellent neighborhood by immediately focusing on the correct reaction (epoxide/anhydride copolymerization) with the correct catalyst family (chromium-salen complexes).. Loser failure mode: Group B's top-ranked papers are significant mismatches, focusing on catalysts with different metals (Mg, Zn, Mn), which dilutes the target's specific context. Although Group B contains a highly relevant paper on dinuclear chromium catalysts at rank 5, its poor ranking and the noise at the top make it a much weaker neighborhood overall..
\item \textbf{A Vygotskian sociocultural perspective on immersion education} (field\_12\_k06\_sum\_minus\_emb\_vs\_consensus\_bank\_n48, $k$=6, over\_regularized\_consensus, gap=3.0). Winner advantage: Group A provides an excellent, highly coherent neighborhood focused on the target's specific topic: the use of L1 in L2 learning from a sociocultural perspective. Its top-ranked papers are directly relevant, and it even includes a paper on the target's specific application context of 'immersion education'.. Loser failure mode: Group B is much broader and more diffuse, focusing on the general sociocultural theory in second language acquisition rather than the target's specific research question, and it completely misses the immersion education context..
\end{itemize}
